\NeedsTeXFormat{LaTeX2e}
\documentclass[a4paper,12pt,
headsepline,        % Linie Kopfzeile / Text
oneside,            % einseitig
bibtotoc,           % Lit.verz. ins Inh.verz.
pointlessnumbers,   % kein Punkt nach letzter Gliederungsziffer
%DIV=15,              % Satzspiegel auf 15er Raster, schmalere Ränder   
BCOR15mm             % mehr li. Rand zum Binden der Arbeit
%,draft
]{scrbook}
\KOMAoptions{DIV=last}

\pagestyle{headings}
\usepackage{blindtext}

% für Texte in deutscher Sprache
\usepackage[ngerman]{babel}
\usepackage[utf8]{inputenc}
\usepackage[T1]{fontenc}

% Helvetica als Standard-Dokumentschrift
\usepackage[scaled]{helvet}
\renewcommand{\familydefault}{\sfdefault} 

\usepackage{graphicx}

% Tabellen mit fester Gesamtbreite und variabler Spaltenbreite
\usepackage{tabularx} 
%\begin{tabularx}{Breite}    {Spaltendefinition} ... \end{tabularx}
%\begin{tabularx}{\textwidth}{XXl} ...               \end{tabularx}

\usepackage{url}              % \url{http://...} in Schreibmaschinenschrift
\usepackage{color}            % zum Setzen farbigen Textes
\usepackage{amssymb, amsmath} % Pakete für Mathe-Umgebungen und -Symbole


\usepackage{setspace} % Paket für div. Abstände, z.B. ZA
\setlength{\parindent}{0pt}
\setlength{\parskip}{1.4ex plus 0.35ex minus 0.3ex}

% Tiefe, bis zu der Überschriften in das Inhaltsverzeichnis kommen
\setcounter{tocdepth}{3} % ist Standard

% Beispiele für Quellcode
\usepackage{listings}
\lstset{language=Java,
  showstringspaces=false,
  frame=single,
  numbers=left,
  basicstyle=\ttfamily,
  numberstyle=\tiny}

% hier Namen etc. einsetzen
% TODO hier anpassen
\newcommand{\fullname}{Florian Krötz}
\newcommand{\email}{florian.kroetz@uni-ulm.de}
\newcommand{\titel}{High Performance implementation of QR decomposition}
\newcommand{\jahr}{2017}
\newcommand{\matnr}{884948}
\newcommand{\gutachterA}{Prof.~Dr.~Streng Geheim}
\newcommand{\gutachterB}{Prof.~Dr.~Un Leserlich}
\newcommand{\betreuer}{Betreuername}

% hier die Fakultät auswählen
%\newcommand{\fakultaet}{---  Im Quellcode anpassen nicht vergessen! ---}
%\newcommand{\fakultaet}{Ingenieurwissenschaften, Informatik und Psychologie}
\newcommand{\fakultaet}{Mathematik und\\Wirtschafts-\\wissenschaften}
%\newcommand{\fakultaet}{Medizin}
%\newcommand{\fakultaet}{Naturwissenschaften}

% hier das Institut einsetzen
\newcommand{\institut}{--- Institut im Quellcode anpassen nicht vergessen! ---}

%color in tables
\usepackage{colortbl}
\definecolor{Gray}{rgb}{0.80784, 0.86667, 0.90196} %dunkelblau
\definecolor{Lightgray}{rgb}{0.9176, 0.95, 0.95686} %hellblau
\definecolor{Akzent}{rgb}{0.6627, 0.63529, 0.55294} %akzentfarbe
\setlength{\arrayrulewidth}{0.1pt}


\pdfinfo{
  /Author (\fullname)
  /Title (\titel)
  /Producer     (pdfeTex 3.14159-1.30.6-2.2)
  /Keywords ()
}

\usepackage{hyperref}
\hypersetup{
pdftitle=\titel,
pdfauthor=\fullname,
pdfsubject={Diplomarbeit},
pdfproducer={pdfeTex 3.14159-1.30.6-2.2},
colorlinks=false,
pdfborder=0 0 0	% keine Box um die Links!
}

%Trennungsregeln
\hyphenation{Sil-ben-trenn-ung}

\begin{document}
\frontmatter

% Titelseite
\thispagestyle{empty}
\begin{addmargin*}[4mm]{-10mm}

\includegraphics[height=1.8cm]{images/unilogo_bild}
\hfill
\includegraphics[height=1.8cm]{images/unilogo_wort}\\[1em]

{\footnotesize
%{\bfseries Universität Ulm} \textbar ~89069 Ulm \textbar ~Germany
\hspace*{115mm}\parbox[t]{35mm}{\bfseries Fakultät für\\
\fakultaet\\

\mdseries \institut}\\[2cm]

\parbox{120mm}{\bfseries \LARGE \titel}\\[2.5em]
{\footnotesize Bachelorarbeit an der Universität Ulm}\\[3em]

{\footnotesize \bfseries Vorgelegt von:}\\
{\footnotesize \fullname\\\email}\\[2em]
{\footnotesize \bfseries Gutachter:}\\                     
{\footnotesize\gutachterA\\
\gutachterB}\\[2em]
{\footnotesize \bfseries Betreuer:}\\ 
{\footnotesize\betreuer}\\\\
{\footnotesize\jahr}
}
\end{addmargin*}


% Impressum
\clearpage
\thispagestyle{empty}
{ \small
  \flushleft
  Fassung \today \\\vfill
  \copyright~\jahr~\fullname\\[0.5em]
% Wenn Sie Ihre Arbeit unter einer freien Lizenz bereitstellen möchten, können Sie die nächste Zeile in Ihren Code aufnehmen. Bitte beachten Sie, dass Sie hierfür an allen Inhalten, inklusive enthaltener Abbildungen, die notwendigen Rechte benötigen! Beim Veröffentlichungsexemplar Ihrer Dissertation achten Sie bitte darauf, dass der Lizenztext nicht den Angaben in den Metadaten der genutzten Publikationsplattform widerspricht. Nähere Information zu den Creative Commons Lizenzen erhalten Sie hier: https://creativecommons.org/licenses/
%This work is licensed under the Creative Commons Attribution 4.0 International (CC BY 4.0) License. To view a copy of this license, visit \href{https://creativecommons.org/licenses/by/4.0/}{https://creativecommons.org/licenses/by/4.0/} or send a letter to Creative Commons, 543 Howard Street, 5th Floor, San Francisco, California, 94105, USA. \\
  Satz: PDF-\LaTeXe
}


% ab hier Zeilenabstand etwas größer 
\setstretch{1.1}

%\onehalfspacing % nur dann, wenn gefordert; ist sehr groß!!

\tableofcontents

\mainmatter
% TODO Kapitel einbinden
\chapter{QR factorisation}

\section{LAPACK}
Mathe \cite{DGEQR2}
\begin{align}
	H &= I - \tau \omega \omega^T \\
	\tau &= \frac{\alpha - \beta}{\beta} \\
	\alpha &= A(i,i)\\
	\beta &= \text{sign}(\alpha) \left|\sqrt{\alpha^2 + \|x\|^2}\right|\\
	x &= A(i+1:m,i)\\
	\omega &= A(i+1:m,i) * \frac{1}{\alpha - \beta}
\end{align}
Algorithmus
\begin{lstlisting}
householderVektor(Vektor v, alpha, tau)
  beta = sign(sqrt(alpha ^2 + norm(x)^2),alpha)
  tau = (alpha - beta) / beta	
  scal(1/(alpha - beta), v)
\end{lstlisting}
\begin{lstlisting}
tau=zeros(min(m,n))
for i = 0 : min(m,n)
  householderVektor(A(i+1:m,i), A(i,i), tau(i)) 
  if (i < n && tau != 0)
    AII = A(i,i)
    A(i,i)= 1
    A = A - tau *w(w'*A) // MV und rank1
    A(i,i) = AII
\end{lstlisting}

\section{NUM3 Urban}
Mathe
\begin{align}
	H &= I - 2 \frac{\omega \omega^T}{\omega^T \omega}\\
	\omega_1 &= \frac{x - \alpha e_1}{x_1 - \alpha}\\
	\alpha ^2 &= \|x\|^2 
\end{align}
Algorithmus
\begin{lstlisting}
householderVektor(Vektor x, omega, beta)
  n = length(x)
  if n> 1
    sigma = x(2:end)'*x(2:end);
    if sigma==0
      beta = 0;
    else
      mu = sqrt(x(1)^2+sigma);
      if x(1)<=0
        tmp = x(1) - mu;
      else
        tmp = -sigma / (x(1) + mu);
      end
      beta = 2*tmp^2/(sigma + tmp^2);
      x(2:end) = x(2:end)/tmp;
    end
    v = [1;x(2:end)];
  else
    beta = 0;
    v = 1;
  end
\end{lstlisting}

\begin{lstlisting}
for i = i:n
  housevector(A(i:m, i), w, beta)
  A(i:m,i:n) = (I(m-i+1) - beta * w * w')*A(i m,i:n)
  if i < m
    A(i + 1 : m, i) = w(2:m-i+1)
\end{lstlisting}



\section{Fragen}
\begin{enumerate}
	\item Was mach ich bei LAPACK mit dem $ \tau $ ?
\end{enumerate}



\appendix
% TODO Anhänge einbinden
\chapter{Quelltexte}

In diesem Anhang sind einige wichtige Quelltexte aufgeführt.

\begin{lstlisting}
public class Hello {
    public static void main(String[] args) {
        System.out.println("Hello World");
    }
}
\end{lstlisting}


\backmatter
\nocite{Knappen2009}
\nocite{Mittelbach2005}
\nocite{Schlosser2014}
\nocite{Sturm2012}
\nocite{Voss2010}

%\bibliographystyle{plaindin} % Nummern und alphabetisch sortiert
%\bibliographystyle{alphadin} % Buchstaben und sortiert
\bibliographystyle{abbrvdin} % Nummern und abgekürzte Namen
%\bibliographystyle{unsrtdin} % Nummern und unsortiert
\bibliography{bibliography}


\clearpage
\thispagestyle{empty}

Name: \fullname \hfill Matrikelnummer: \matnr \vspace{2cm}

\minisec{Erklärung}

Ich erkläre, dass ich die Arbeit selbständig verfasst und keine anderen als die angegebenen Quellen und Hilfsmittel verwendet habe.\vspace{2cm}

Ulm, den \dotfill

\hspace{10cm} {\footnotesize \fullname}
\end{document}
